\section{Introduction}

\subsection{The Fermi Paradox and Its Partial Solutions}

The Fermi Paradox---the absence of observational evidence for extraterrestrial civilizations despite their apparent high probability \cite{hart1975}---admits two classes of solution. \emph{Filtering} solutions posit a barrier that destroys civilizations or prevents their emergence, explaining the complete absence of detectable signals. \emph{Stagnation} solutions posit that civilizations persist but do not expand or emit detectable signatures, explaining the absence of interstellar presence but not necessarily the absence of all technosignatures.

Hanson's Great Filter hypothesis \cite{hanson1998} belongs to the first class: a critical barrier---behind us or ahead of us---that virtually no civilization survives. Proposed filters include nuclear annihilation, ecological collapse, and artificial intelligence misalignment \cite{bostrom2002,bostrom2014}. These explanations share a focus on exogenous or technological threats.

We argue that both classes of solution are necessary, and that the most dangerous filter is \emph{endogenous}. Civilizations do not need an external catastrophe to fail. The internal dynamics of social hierarchy---specifically, the tendency of elites to optimize resource extraction over systemic adaptation---can dissipate adaptive potential from within, trapping a civilization in a stable low-technology equilibrium from which interstellar expansion is impossible. This is the \emph{Great Plateau}. At a sufficiently advanced technological threshold, this stagnation may transition to \emph{filtration}: the elite class no longer requires a human population base, and the civilization's technosignatures fade below the threshold of astronomical detectability.

\subsection{Intraspecific Economic Cannibalism}

We formalize this mechanism through the concept of \emph{intraspecific economic cannibalism}: the process by which elites, competing for guaranteed access to resources, extract societal surplus without productive reinvestment. This is not a normative claim about any particular political system. It is a structural tendency present in any sufficiently complex society where some actors can capture disproportionate resources and where the institutions that might restrain them are themselves vulnerable to capture. The mathematical theory of complex societal collapse has been developed by \citet{tainter1988,turchin2003,turchin2016}, who identified mechanisms including diminishing returns on complexity, elite overproduction, and fiscal fragility. Our contribution extends this tradition by formalizing the interaction between extraction, technology, and cooperation in a dynamical system and connecting it to the astrophysical problem of the Fermi Paradox.

The mechanism operates through three interdependent variables: technology ($T$), the extraction rate ($K$), and cooperative capacity ($C$). When extraction exceeds a critical threshold $K_{\text{crit}}(T)$, the feedback loops that sustain technological development reverse, and the system enters the feudal trap attractor. The model makes no distinction between political systems---it describes a universal tension between rent-seeking and cooperation.

\subsection{From Plateau to Filter}

A civilization trapped at the feudal plateau retains basic technology ($T \approx 0.1$), population, and---crucially---detectable technosignatures. It has not passed through a Great Filter; it has stalled on a Great Plateau. The transition from plateau to filter requires an additional mechanism.

The EDAP model identifies a candidate mechanism in the technofeudal threshold $T_{\text{sing}}$. Below this threshold, elites depend on human populations for labor, consumption, and taxation---even in highly extractive regimes, there is a lower bound on tolerable population decline. Beyond $T_{\text{sing}}$, automation, artificial intelligence, and robotic production decouple elite welfare from population size. The extraction ceiling $K_{\max}(T) \rightarrow 1.0$---there is no physical limit to extraction when the means of production are autonomous. Shadow cooperation, the last defense against total extraction, is suppressed by pervasive surveillance technologies. A civilization crossing $T_{\text{sing}}$ with $K > K_{\text{crit}}$ may undergo \emph{filtration}: population collapse, technological simplification below the threshold of astronomical detectability, and permanent disappearance from the observable universe.

\textbf{We emphasize that this filter transition is a qualitative implication of the model's mathematical structure, not a quantitative prediction verified by simulation.} The current model does not track population, technosignature strength, or astronomical detectability. Establishing the plateau-to-filter transition as a robust, calibrated mechanism requires explicit modeling of these variables, which we defer to future work.

This paper formalizes the plateau mechanism mathematically. The filter mechanism follows as a qualitative consequence of the model's structure at the technofeudal threshold. Together, they provide a unified endogenous framework for addressing the Fermi Paradox.
